\exercisegroup

\begin{enumerate}
\def\labelenumi{\arabic{enumi})}
\tightlist
\item
  对于每个 \(\theta \in \mathbf{R}\)，令 \(C_{\theta}\) 为平面中端点为 \((0,0)\) 和 \((\cos(2\pi\theta), \sin(2\pi\theta))\) 的线段；置
\end{enumerate}

\[
\mathrm{C} = \bigcup_ {\theta \in \mathbf {Q} \cap [ 1 / 6, 1 / 4 ]} \mathrm{C} _ {\theta}.
\]

\begin{enumerate}
\def\labelenumi{\alph{enumi})}
\item
  空间 C 在 0 处可缩，但在任何其他点处都不可缩。
\item
  令 D 为 C 的平移 \((0, n) + \mathrm{C}\)（\(n \in Z\)）之并。证明空间 D 同伦等价于一点，但在其任何一点处都不可缩。
\end{enumerate}

\begin{enumerate}
\def\labelenumi{\arabic{enumi})}
\setcounter{enumi}{1}
\tightlist
\item
  设 B 是 \(\mathbf{R}^{n}\) 的子集，\(a\) 是 B 内部的一点；记 S 为 B 的边界。假设 B 关于 \(a\) 是星形的，且每条起点为 \(a\) 的半直线都与 S 交于唯一一点。令 \(f\colon\mathbf{R}_{+}^{*}\times\mathrm{S}\to\mathbf{R}^{n}-\{a\}\) 为由 \(f(t,y)=ty+(1-t)a\) 定义的映射。证明 \(f\) 是同胚，并且满足 \(f([0,1]\times\mathrm{S})=\mathrm{B}-\{a\}\) 和 \(f(\{1\}\times\mathrm{S})=\mathrm{S}\)。
\end{enumerate}

当 B 是 \(R^{n}\) 的凸、闭且有界子集，a 是 B 的附着点时，该结果尤其适用（参见 I，第 123 页，命题 2 以及 EVT，II，第 15 页，命题 16）。

\begin{enumerate}
\def\labelenumi{\arabic{enumi})}
\setcounter{enumi}{2}
\item
  设 X 是将空间 R 中的子集 Q 缩塌所得的拓扑空间（III，第 252 页，定义 10）。典范满射 \(p: R \rightarrow X\) 既非开映射也非闭映射。
\item
  设 X 是 R 中的区间 \([-1,1]\)，令 A = X - \{0\}。X 到 X/A 的典范映射是同伦等价，但不存在起点为 \(\mathrm{Id}_{X}\) 且满足 III 第 252 页命题 10 的假设的同伦 σ。
\item
  令 A 为 \(\mathbf{R}\) 中形如 \(1/n\) 的实数集合的闭包，其中 \(n\) 遍历满足 \(\geqslant 1\) 的整数集合。令 \(i\) 为 A 到 \(\mathbf{R}\) 的典范嵌入。由 \(\mathrm{Côn}(i)\) 到 \(\mathbf{R}/\mathrm{A}\) 的典范映射不是同伦等价。
\item
  设 X 是可度量化的可数型拓扑空间。称 X 为绝对收缩核\(^{(1)}\)，如果对于每个可度量化的可数型拓扑空间 Y、Y 的每个闭子集 B 以及每个连续映射 \(f: B \to X\)，映射 f 都能延拓为从 Y 到 X 的连续映射。
\end{enumerate}

\begin{enumerate}
\def\labelenumi{\alph{enumi})}
\item
  设 X 是可度量化的可数型拓扑空间。要使 X 成为绝对收缩核，必要且充分条件是：对于每个可度量化的可数型空间 Y 以及每个从 X 到 Y 的连续、单射且闭映射 f，都存在连续映射 \(g: Y \to X\)，使得 \(g \circ f = Id_{X}\)。
\item
  证明空间 I 和 R 是绝对收缩核。
\item
  证明绝对收缩核的可数族的乘积空间是绝对收缩核。
\item
  设 X 是绝对收缩核，A 是 X 的子空间且是 X 的收缩核。证明 A 是绝对收缩核。
\end{enumerate}

\begin{enumerate}
\def\labelenumi{\arabic{enumi})}
\setcounter{enumi}{6}
\tightlist
\item
  设 X 是可度量化的可数型拓扑空间。如果对于每个可度量化的可数型拓扑空间 Y、Y 的每个闭子集 B 以及每个连续映射 \(f: B \rightarrow X\)，映射 f 都能延拓为从 B 的某个邻域到 X 的连续映射，则称 X 为绝对邻域收缩核。
\end{enumerate}

\begin{enumerate}
\def\labelenumi{\alph{enumi})}
\item
  设 X 是可度量化的可数型拓扑空间。要使 X 成为绝对邻域收缩核，必要且充分条件是：对于每个可度量化的可数型空间 Y 以及每个从 X 到 Y 的连续、单射且闭映射 f，都存在 \(f(\mathrm{X})\) 的一个邻域 U 以及连续映射 \(g: U \to X\)，使得 \(g \circ f = Id_{X}\)。
\item
  设 X 是绝对邻域收缩核，A 是 X 的子集。若 A 有一个邻域 U，且 A 是 U 的收缩核，则 A 是绝对邻域收缩核。特别地，X 的每个开子集以及 X 的每个收缩核都是绝对邻域收缩核。
\item
  证明绝对邻域收缩核的有限族的乘积空间是绝对邻域收缩核。
\end{enumerate}

d) 证明球面 \(\mathbf{S}_n\) 是绝对邻域收缩核。

\begin{enumerate}
\def\labelenumi{\alph{enumi})}
\setcounter{enumi}{4}
\item
  设 X 是可度量化的可数型拓扑空间，\(A_{1}\) 和 \(A_{2}\) 是 X 的闭子集，且 \(X = A_{1} \cup A_{2}\)。假设 \(A_{1} \cap A_{2}\) 是绝对邻域收缩核。要使 X 成为绝对邻域收缩核，必要且充分条件是 \(A_{1}\) 和 \(A_{2}\) 也具有同样性质。
\item
  设 X 是可度量化的可数型拓扑空间。若 X 是两个绝对邻域收缩核开子集之并，则 X 是绝对邻域收缩核。
\item
  设 X 是可度量化的可数型拓扑空间，且每一点都有一个绝对邻域收缩核的邻域。证明 X 是绝对邻域收缩核。
\end{enumerate}

\begin{enumerate}
\def\labelenumi{\arabic{enumi})}
\setcounter{enumi}{7}
\tightlist
\item
  设 X 是可度量化的拓扑空间；设 d 是定义其拓扑的度量。令 \(\mathcal{C}_{b}(\mathrm{X};\mathbf{R})\) 为 X 上连续且有界的实值函数构成的向量空间；给它赋予由 \(\|f\| = \sup_{x \in X} |f(x)|\)（\(f \in \mathcal{C}_{b}(\mathrm{X};\mathbf{R})\)）定义的范数。
\end{enumerate}

\begin{enumerate}
\def\labelenumi{\alph{enumi})}
\item
  定义映射 \(j\colon X \to \mathcal{C}_{b}(X; \mathbf{R})\)，其定义为 \(j(x)\colon y \mapsto d(x, y)/(1 + d(x, y))\)。证明映射 j 诱导从 X 到其像的同胚。
\item
  令 V 为 \(\mathcal{C}_{b}(\mathrm{X};\mathbf{R})\) 中所有包含 \(j(\mathrm{X})\) 的凸子集之交。证明 \(j(\mathrm{X})\) 在 V 中是闭的。
\item
  设 X 是可度量化的可数型空间。证明存在 \(I^{N}\) 的凸子集 Z，使得 X 同胚于 Z 的一个闭子集。
\end{enumerate}

\begin{enumerate}
\def\labelenumi{\arabic{enumi})}
\setcounter{enumi}{8}
\tightlist
\item
  如果 X、Y 是拓扑空间，U 是 X 的子集族，则称映射 \(f, g: Y \to X\) 为 U-邻近的，如果对于每个 \(y \in Y\)，存在 \(U \in U\)，使得 \(f(y)\) 和 \(g(y)\) 都属于 U；称同伦 \(\sigma: Y \times I \to X\) 为 U-小的，如果对于每个 \(y \in Y\)，存在 \(U \in U\)，使得 \(\sigma(y, t) \in U\) 对所有 \(t \in I\) 都成立。
\end{enumerate}

设 X 是 \(I^{N}\) 的凸子集 Z 的闭子空间（参见习题 8）；假设 X 是绝对邻域收缩核。

\begin{enumerate}
\def\labelenumi{\alph{enumi})}
\item
  设 U 是由 \(I^{N}\) 的开子集构成的 X 的覆盖。证明存在 \(I^{N}\) 的开凸子集族 V，使得对于每个开集 \(v \in V\)，存在开集 \(u \in U\)，满足 \(v \cap X \subset u\)，并且存在 X 到族 V 的并集中的典范嵌入的连续收缩映射 r。
\item
  设 Y 是拓扑空间，B 是 Y 的闭子空间。设 \(f, g: \mathrm{Y} \to \mathrm{X}\) 是 \(\mathcal{V}\)-邻近的连续映射，\(\sigma: \mathrm{B} \times \mathbf{I} \to \mathrm{X}\) 是 \(\mathcal{V}\)-小的同伦，且其起点等于 \(f|_{\mathrm{B}}\)。
\end{enumerate}

  证明存在 Y 中 B 的一个邻域 W，以及从子空间 \((\mathbf{Y} \times \{0,1\}) \cup (\mathbf{W} \times \mathbf{I})\)（它是 \(Y \times I\) 的子空间）到 X 的连续映射 h，使得 \(h(y,0) = f(y)\) 和 \(h(y,1) = g(y)\) 对所有 \(y \in Y\) 都成立，并且 \(h(y,t) = \sigma(y,t)\) 对所有 \(y \in B\) 和每个 \(t \in I\) 都成立。证明存在包含于 W 的 B 的一个邻域 \(W'\)，使得 \(h \mid W' \times I\) 是 \(\mathcal{V}\)-小同伦。

\begin{enumerate}
\def\labelenumi{\alph{enumi})}
\setcounter{enumi}{2}
\tightlist
\item
  设 \(u \colon \mathrm{Y} \to \mathbf{I}\) 是连续映射，满足 \(u(y) = 1\) 对 \(y \in \mathrm{B}\) 成立，且 \(u(y) = 0\) 在 \(\mathrm{Y} - \mathrm{W}\) 的某个邻域上成立；置
\end{enumerate}

\[
h ^ {\prime} (y, t) = \left\{ \begin{array}{l l} (1 - u (y)) ((1 - t) f (y) + t g (y)) + u (y) h (y, t) & \text { if } y \in \mathrm{V} \\ (1 - t) f (y) + t g (y) & \text { otherwise. } \end{array} \right.
\]

证明 \(h'\) 连续且其像包含于 V。

\(d)\) 推出存在同伦 \(\widetilde{\sigma}\colon\mathbf{Y}\times\mathbf{I}\to\mathbf{X}\)，其起点为 \(f\)、终点为 \(g\)，且是 \(\mathcal{U}\)-小的。

\begin{enumerate}
\def\labelenumi{\arabic{enumi})}
\setcounter{enumi}{9}
\tightlist
\item
  设 X 是拓扑空间。假设 X 有一个开覆盖 U，使得对于每个拓扑空间 Y 和 Y 的每个闭子空间 B，从 Y 到 X 的两个 U-邻近映射 \(f_{0}, f_{1}\)，只要存在连接 \(f_{0} \mid B\) 与 \(f_{1} \mid B\) 的 U-小同伦，它们就是同伦的。
\end{enumerate}

\begin{enumerate}
\def\labelenumi{\alph{enumi})}
\item
  设 a 是 X 的一点，U 是 U 中满足 \(a \in U\) 的一个元素。证明存在同伦 \(\sigma: U \times I \to X\)，使得 \(\sigma(x,0) = a\)、\(\sigma(x,1) = x\) 且 \(\sigma(a,t) = a\) 对所有 \(x \in X\) 和每个 \(t \in I\) 成立。
\item
  设 V 是 a 的开邻域，满足 \(\sigma(\mathrm{V} \times \mathbf{I}) \subset \mathrm{U}\)。证明 V 是绝对邻域收缩核。
\item
  证明 X 是绝对邻域收缩核。
\end{enumerate}

\begin{enumerate}
\def\labelenumi{\arabic{enumi})}
\setcounter{enumi}{10}
\tightlist
\item
  设 X 是绝对邻域收缩核。
\end{enumerate}

\begin{enumerate}
\def\labelenumi{\alph{enumi})}
\item
  证明 X 是“半局部可缩”的，即 X 的每一点 x 都有一个邻域 V，使得 (V, x) 到 X 的典范嵌入严格同伦于像为 x 的常值映射。（使用习题 9。）
\item
  设 A 是紧拓扑空间。证明空间 \(\mathcal{C}_{c}(\mathrm{A};\mathrm{X})\) 的道路连通分支是开的。
\item
  证明 X 局部道路连通。
\end{enumerate}

\begin{enumerate}
\def\labelenumi{\arabic{enumi})}
\setcounter{enumi}{11}
\tightlist
\item
  设 X 和 Y 是拓扑空间；如果对于每个连续映射 F: \(Y \rightarrow X\)、Y 的每个闭子空间 B 以及每个同伦 \(\sigma\colon\mathrm{B}\times\mathbf{I}\to\mathrm{X}\)（其起点为 F \textbar{} B），都存在一个同伦 \(\widetilde{\sigma}\colon\mathrm{Y}\times\mathbf{I}\to\mathrm{X}\)，其起点为 F，并延拓 \(\sigma\)，则称连续映射的同伦扩张性质对目标 X 和源 Y 成立。
\end{enumerate}

\begin{enumerate}
\def\labelenumi{\alph{enumi})}
\item
  设 X 是绝对邻域收缩核的拓扑空间。证明连续映射的同伦扩张性质对目标 X 和可度量化的可数型空间作为源成立。
\item
  证明下列两个性质等价：
\end{enumerate}

\begin{enumerate}
\def\labelenumi{(\roman{enumi})}
\item
  空间 X 是绝对邻域收缩核；
\item
  X 的每一点都有一个邻域 V，使得对于每个可度量化的可数型拓扑空间 Y、Y 的每个闭子空间 B 以及从 B 到 V 的每个连续映射，都能延拓为从 Y 到 X 的连续映射。
\end{enumerate}

\begin{enumerate}
\def\labelenumi{\alph{enumi})}
\setcounter{enumi}{2}
\item
  假设 X 是半局部可缩的（参见习题 9），且连续映射的同伦扩张性质对目标 X 成立。证明 X 是绝对邻域收缩核。\(^{(2)}\)
\item
  证明空间 X = Q 不是绝对邻域收缩核，但连续映射的同伦扩张性质对目标 X 成立。
\end{enumerate}

\begin{enumerate}
\def\labelenumi{\arabic{enumi})}
\setcounter{enumi}{12}
\item
  设 X 是绝对邻域收缩核的拓扑空间。设 A 是 X 的子集，使得 A 到 X 的典范嵌入是同伦等价和余纤维化。存在 X 到 A 的收缩。
\item
  设 X 是拓扑空间 I，置 A = \{0,1\}。证明存在连接 X 的恒等映射与常值映射的同伦，但典范映射 X → X/A 不是同伦等价。
\item
  设 X 是 \([0,1]\times[0,1]\) 的子空间，由点对 \((x,y)\) 构成，其中满足 \(x\in\{0,1\}\)、或 \(y\in\{0,1\}\)、或 \(x\neq0\) 且 \(1/x\in N\)。令 A = \(\{0\}\times[0,1]\)。证明从 X 到 X/A 的典范映射不是同伦等价，尽管 A 可缩。
\item
  设 X 和 Y 是拓扑空间，\(f: X \to Y\) 是连续映射。令 \(\varphi: C(X) \to \text{Côn}(f)\) 为典范映射。证明空间 \(Y/f(X)\) 与 \(\text{Côn}(f)/\varphi(C(X))\) 同胚。
\item
  设 X 是拓扑空间，A 是 X 的子空间。记 \(i\colon A \to X\) 和 \(j\colon \operatorname{Cyl}(i) \to X \times I\) 为典范嵌入。
\end{enumerate}

\begin{enumerate}
\def\labelenumi{\alph{enumi})}
\item
  假设 X = R 且 \(A = ]0, 1[\)；证明映射 j 不是严格映射。
\item
  假设 X 的每一点都有可数的闭邻域基。映射 j 严格的必要且充分条件是 A 闭。
\item
  假设 A 是 Alexandroff 半直线（TG，IV，第 49 页，习题 12），X 是它的 Alexandroff 紧化。证明映射 j 是严格映射，尽管 A 在 X 中不是闭的。
\end{enumerate}

\begin{enumerate}
\def\labelenumi{\arabic{enumi})}
\setcounter{enumi}{17}
\tightlist
\item
如果从 X 到 X × X 的对角映射是余纤维化，则称拓扑空间 X 局部等连通。
\end{enumerate}

\begin{enumerate}
\def\labelenumi{\alph{enumi})}
\tightlist
\item
  证明局部等连通空间的有限族的乘积是局部等连通的。
\end{enumerate}

设 X 是局部等连通空间。

\begin{enumerate}
\def\labelenumi{\alph{enumi})}
\setcounter{enumi}{1}
\item
  证明存在由 X 的开子集序列 \((\mathrm{U}_{n})_{n\in\mathbf{N}}\) 构成的 X 的覆盖，使得对于每个 n，\(U_{n}\) 到 X 的嵌入同伦于常值映射。
\item
  证明 X 是豪斯多夫空间。
\item
  设 \(u: X \to I\) 是连续映射，令 U 为所有满足 \(x \in X\) 且 \(u(x) > 0\) 的点 x 所成的集合。证明 U 的道路连通分支是开的。特别地，X 的道路连通分支是开的。
\item
  设 A 是 X 的闭子空间，且是 X 的收缩核。证明 A 局部等连通，且对 (X, A) 同伦扩张性质成立。（“Dyer–Eilenberg 定理”。）
\item
  设 A 是 X 的子空间，且对 (X, A) 同伦扩张性质成立。证明 A 局部等连通。
\end{enumerate}

\begin{enumerate}
\def\labelenumi{\arabic{enumi})}
\setcounter{enumi}{18}
\tightlist
\item
  设 B 是拓扑空间，\(u \colon \mathrm{X} \to \mathbf{I}\) 是连续映射，且 \(\mathrm{A} = u^{-1}(0)\)。
\end{enumerate}

\begin{enumerate}
\def\labelenumi{\alph{enumi})}
\item
  证明从 \(X \times I\) 到自身的映射 p（由 \(p(x,t) = (x,tu(x))\) 给出）是闭映射。
\item
  设 Y 是拓扑空间，\(\sigma\colon\mathrm{X}\times\mathbf{I}\to\mathrm{Y}\) 是在 A 上固定的同伦。令 \(\tau\colon\mathrm{X}\times\mathbf{I}\to\mathrm{Y}\) 为由 \(\tau(x,t)=\sigma(x,\inf(1,t/u(x)))\)（当 \(u(x)>0\) 时）以及 \(\tau(x,t)=\sigma(x,1)\)（否则）给出的映射。证明 \(\tau\) 连续。
\item
  设 C 是 X 的子空间；假设对 (X, C) 同伦扩张性质成立。设 \(f \colon \mathrm{X} \to \mathrm{Y}\) 是连续映射，\(\sigma \colon \mathrm{C} \times \mathbf{I} \to \mathrm{Y}\) 是起点为 \(f \mid \mathrm{C}\) 且在 \(\mathrm{A} \cap \mathrm{C}\) 上固定的同伦。证明存在同伦 \(\tau \colon \mathrm{X} \times \mathbf{I} \to \mathrm{Y}\)，其起点为 \(f\)，延拓 \(\sigma\)，且在 \(\mathrm{A}\) 上固定。
\end{enumerate}

\begin{enumerate}
\def\labelenumi{\arabic{enumi})}
\setcounter{enumi}{19}
\tightlist
\item
  设 X 是拓扑空间，A、B、C 是 X 的闭子空间，且 \(A = B \cap C\)。
\end{enumerate}

\begin{enumerate}
\def\labelenumi{\alph{enumi})}
\item
  若对 (X, B)、(X, C) 和 (X, A) 同伦扩张性质成立，证明对 (X, B ∪ C) 也成立。
\item
  假设对 (X, A) 和 (X, B ∪ C) 同伦扩张性质成立。证明对 (X, B) 和 (X, C) 也成立。
\item
  假设 X 是正规空间，且 \(A \subset \mathring{B} \cup \mathring{C}\)。若对 (X, B) 和 (X, C) 同伦扩张性质成立，证明对 (X, A) 也成立。
\end{enumerate}

\begin{enumerate}
\def\labelenumi{\arabic{enumi})}
\setcounter{enumi}{20}
\tightlist
\item
  设 X 是仿紧拓扑空间，\(\mathcal{A}\) 是 X 的闭子集所成的集合，并且每一对元素 B、C 都满足 \(B \cap C \in \mathcal{A}\)，其中 B、C 属于 \(\mathcal{A}\)。令 \(A = \bigcup_{C \in \mathcal{A}} C\)；假设 X 的每一点 x 都有一个 x 的邻域 V 以及一个有限子集 \(\mathcal{A}_x\)（它包含于 \(\mathcal{A}\)），使得 \(V \cap A = V \cap \bigcup_{C \in \mathcal{A}_x} C\)。a) 证明 A 是 X 的闭子集。
\end{enumerate}

\begin{enumerate}
\def\labelenumi{\alph{enumi})}
\setcounter{enumi}{1}
\tightlist
\item
  假设对于所有 \(B \in A\)，对 (X, B) 同伦扩张性质成立。证明对 (X, A) 也成立。
\end{enumerate}

\begin{enumerate}
\def\labelenumi{\arabic{enumi})}
\setcounter{enumi}{21}
\tightlist
\item
  设 X 和 Y 是拓扑空间，A 是 X 的子空间，B 是 Y 的子空间。
\end{enumerate}

\begin{enumerate}
\def\labelenumi{\alph{enumi})}
\item
  假设 (X,A) 和 (Y,B) 同伦等价。若对 (X,A) 同伦扩张性质成立，证明对 (Y,B) 也成立。
\item
  反之，假设对 (X, A) 和 (Y, B) 同伦扩张性质成立。设 \(f \colon X \to Y\) 是同伦等价，并且通过限制到子空间诱导出从 A 到 B 的同胚。证明 \(f\) 是从对 (X, A) 到对 (Y, B) 的同伦等价。
\end{enumerate}

\begin{enumerate}
\def\labelenumi{\arabic{enumi})}
\setcounter{enumi}{22}
\tightlist
\item
  设 X 和 Y 是拓扑空间。同伦 \(\sigma\colon\mathrm{X}\times\mathbf{I}\to\mathrm{Y}\) 若存在实数 \(\delta>0\)，使得 \(\sigma(x,t)=\sigma(x,0)\) 对所有 \(x\in\mathrm{X}\) 和每个 \(t\in[0,\delta]\) 都成立，则称其为平稳的。设 A 是 X 的闭子空间；记 \(j\) 为 A 到 X 的典范嵌入。如果对于每个拓扑空间 Y、每个连续映射 \(f\colon\mathrm{X}\to\mathrm{Y}\) 以及每个平稳同伦 \(\sigma\colon\mathrm{A}\times\mathbf{I}\to\mathrm{Y}\)（其起点等于 \(f\mid\mathrm{A}\)），都存在同伦 \(\widetilde{\sigma}\colon\mathrm{X}\times\mathbf{I}\to\mathrm{Y}\)，其起点为 \(f\)，并延拓 \(\sigma\)，则称对 (X,A) 的平稳同伦同伦扩张性质成立。
\end{enumerate}

\begin{enumerate}
\def\labelenumi{\alph{enumi})}
\tightlist
\item
  证明下列性质等价：
\end{enumerate}

\begin{enumerate}
\def\labelenumi{(\roman{enumi})}
\item
  对 (X, A) 平稳同伦扩张性质成立。
\item
  (X, A) 与 \((\mathrm{Cyl}(j), \alpha_{j}(\mathrm{A} \times \{0\}))\) 同伦等价；
\item
  存在连续映射 \(\varphi\colon X\to\mathbf{I}\)，使得 \(\varphi(a)=0\) 对所有 \(a\in A\) 成立；并且存在同伦 \(\sigma\colon X\times\mathbf{I}\to X\)，使得 \(\sigma(x,0)=x\) 对 \(x\in X\) 成立，\(\sigma(a,t)=a\) 对 \((a,t)\in A\times\mathbf{I}\) 成立，并且 \(\sigma(x,1)\in A\)（若 \(\varphi(x)<1\)）。
\item
  存在连续映射 \(\varphi\colon\mathrm{X}\to\mathbf{I}\)、X 的子空间 V 以及同伦 \(\sigma\colon\mathrm{V}\times\mathbf{I}\to\mathrm{X}\)，使得 \(\varphi(x)=1\)（若 \(x\notin\mathrm{V}\)），\(\varphi(a)=0\)（对 \(a\in\mathrm{A}\)），\(\sigma(x,0)=x\)（对 \(x\in\mathrm{V}\)），\(\sigma(a,t)=a\)（对 \((a,t)\in\mathrm{A}\times\mathbf{I}\)），并且 \(\sigma(x,1)\in\mathrm{A}\)（若 \(\varphi(x)<1\)）。
\end{enumerate}

\begin{enumerate}
\def\labelenumi{\alph{enumi})}
\setcounter{enumi}{1}
\tightlist
\item
  假设对 (X, A) 平稳同伦延拓性质成立，并且存在连续函数 \(u \colon X \to I\)，使得 \(u^{-1}(0) = A\)。证明对 (X, A) 同伦扩张性质成立。
\end{enumerate}

\begin{enumerate}
\def\labelenumi{\arabic{enumi})}
\setcounter{enumi}{23}
\tightlist
\item
  设 M 是不可数集，X 是空间 \(I^{M}\)，A 是 X 的子空间 \(\{0\}^{M}\)。
\end{enumerate}

\begin{enumerate}
\def\labelenumi{\alph{enumi})}
\item
  证明对 (X, A) 平稳同伦扩张性质成立（习题 23）。
\item
  证明对 (X, A) 同伦扩张性质不成立。（证明集合 A 不是可数族开集的交。）
\item
  证明对 \((\mathbf{X} \times \mathbf{I}, \mathbf{X} \times \{0\} \cup \mathbf{A} \times \mathbf{I})\) 平稳同伦延拓性质不成立。
\end{enumerate}

\begin{enumerate}
\def\labelenumi{\arabic{enumi})}
\setcounter{enumi}{24}
\tightlist
\item
  \begin{enumerate}
  \def\labelenumii{\alph{enumii})}
  \tightlist
  \item
  设 X 和 Y 是拓扑空间，A 是 X 的闭子空间，B 是 Y 的闭子空间。设 \(f \colon \mathrm{X} \to \mathrm{Y}\) 是连续映射；假设 \(f\) 通过限制到子空间诱导出从 A 到 B 的同胚。若 \(f\) 有同伦左逆，且对 (Y, B) 平稳同伦扩张性质成立（习题 23），则对 (X, A) 也成立。
  \end{enumerate}
\end{enumerate}

\begin{enumerate}
\def\labelenumi{\alph{enumi})}
\setcounter{enumi}{1}
\item
  若 (X, A) 和 (Y, B) 同伦等价，则对 (X, A) 平稳同伦扩张性质成立，当且仅当对 (Y, B) 也成立。
\item
  令 C 为平面中形如 \((t/n,t)\) 的点的集合，其中 \(t \in I\) 且 \(n \in N^{*}\)；令 X 为其闭包，并令 \(A = \{0\} \times I\)。证明对 \((X,A)\) 平稳同伦延拓性质不成立。
\end{enumerate}

\begin{enumerate}
\def\labelenumi{\arabic{enumi})}
\setcounter{enumi}{25}
\item
  给集合 \(X = \{0,1\}\) 赋予拓扑，使开集为 \(\varnothing\)、\(\{0\}\) 和 X。置 \(A = \{0\}\)。证明对 \((X,A)\) 同伦扩张性质成立，但对 \((X \times X,(X \times A) \cup (A \times X))\) 不成立。
\item
  设 X 是拓扑空间，A 是 X 的子空间（不必闭）；记子空间 \((\mathrm{X} \times \{0\}) \cup (\mathrm{A} \times \mathbf{I})\) 为 C，它是 \(X \times I\) 的子空间。假设 C 是 \(X \times I\) 的收缩核。设 U 是 C 的子集，使得 \(\mathrm{U} \cap (\mathrm{X} \times \{0\})\) 和 \(\mathrm{U} \cap (\mathrm{A} \times \mathbf{I})\) 分别是 \(X \times \{0\}\) 和 \(A \times I\) 中的开集。a) 令 \(V_{0}\) 为满足 \(x \in X\) 且 \((x, 0) \in \mathrm{U}\) 的 x 的集合；对于 \(n \geqslant 1\)，令 \(V_{n}\) 为 X 中所有开集 W 的并，其中 \((\mathrm{W} \cap \mathrm{A}) \times [0, 1/n]\) 包含于 U。
\end{enumerate}

证明 \(A \cap V_{0} = \bigcup_{n \geqslant 1} A \cap V_{n}\)；还要证明，X 中满足 \(W \cap A \subset V_{n}\) 的每个开集 W 都包含于 \(V_{n}\)。

\begin{enumerate}
\def\labelenumi{\alph{enumi})}
\setcounter{enumi}{1}
\item
  证明 \(\mathrm{V} \subset \bigcup_{n \geqslant 1} \mathrm{V}_n\)。
\item
  证明 U 在 C 中是开的。
\item
  证明对 (X, A) 同伦扩张性质成立。
\end{enumerate}

\begin{enumerate}
\def\labelenumi{\arabic{enumi})}
\setcounter{enumi}{27}
\item
  设 \(X\) 是拓扑空间，\(A\) 是 X 的子空间，且对 \((X, A)\) 同伦扩张性质成立。证明对 \((X, \overline{A})\) 也成立。
\item
  设 B 是拓扑空间。B 的覆盖 \((\mathrm{V}_{i})_{i\in\mathrm{I}}\) 称为可数化的，如果 B 上存在局部有限的连续单位分解 \((g_{j})_{j\in\mathrm{J}}\)，使得对于每个 \(j\in J\)，\(g_{j}^{-1}(]0,1])\) 包含于某个 \(V_{i}\) 中。
\end{enumerate}

\begin{enumerate}
\def\labelenumi{\alph{enumi})}
\item
  证明 B 是仿紧的，当且仅当 B 的每个开覆盖都是可数化的。
\item
  证明 B 是正规的，当且仅当 B 的每个局部有限开覆盖都是可数化的。
\item
  设 \((\mathrm{V}_{i})\) 是 B 的可数化覆盖，\(B'\) 是拓扑空间，且 \(f: B' \to B\) 是连续映射。证明覆盖 \((f^{-1}(\mathrm{V}_{i}))\)（即 \(B'\) 的覆盖）是可数化的。
\end{enumerate}

\begin{enumerate}
\def\labelenumi{\arabic{enumi})}
\setcounter{enumi}{29}
\item
  设 B 是仿紧拓扑空间，F 是可缩拓扑空间，E 是 B 上纤维为 F 的局部平凡纤维丛。证明，E 的连续截面的芽层是软层。
\item
  设 B 是拓扑空间。设 A 和 V 是 B 的子空间；若存在连续映射 \(\varphi\colon\mathbf{B}\to\mathbf{I}\)，使其在 A 的每一点取值为 1，在 CV 的每一点取值为 0，则称 V 是 A 的一个光环。若对于 B 的每个子空间 A、A 的每个光环 V 以及 E|V 的每个截面 s，都存在 E 的截面 s' 使得 s' \textbar{} A = s \textbar{} A，则称 B-空间 (E,p) 具有截面延拓性质。
\end{enumerate}

\begin{enumerate}
\def\labelenumi{\alph{enumi})}
\item
  设 \((\mathrm{E},p)\) 和 \((\mathrm{E}^{\prime},p^{\prime})\) 是 B-空间；假设存在 B-空间的态射 \(f: E \to E'\) 和 \(g: E' \to E\)，以及一个同伦 \(\sigma: E \times I \to E\)，其起点为 \(Id_{E}\)，终点为 \(g \circ f\)，并且 \(p(\sigma(x,t)) = p(x)\) 对所有 \((x,t) \in E \times I\) 都成立。若 B-空间 \(E'\) 具有截面延拓性质，则 B-空间 E 也具有该性质。
\item
  设 E 是 B-空间，A 是 B 的子空间且是 B 的收缩核。假设诱导的 A-空间 \(E_{A}\) 具有截面延拓性质；证明 B-空间 E 也具有同样的性质。
\item
 设 E 是具有截面延拓性质的 B-空间。令 \(f \colon \mathrm{B} \to \mathbf{I}\) 为连续映射，并令 \(V = f^{-1}(]0,1])\)。证明 V-空间 \(E_V\) 具有截面延拓性质。（令 A 为 V 的子空间，W 为 V 中 A 的一个光环；构造 B 的子空间序列 \((A_n)\) 和 \((V_n)\)，使其满足以下性质：对每个 \(n\)，\(A_n\) 包含于 W 中，\(V_n\) 是 \(A_n\) 的一个光环，且所有 \(A_n\) 的交等于 A。）
\end{enumerate}

\begin{enumerate}
\def\labelenumi{\arabic{enumi})}
\setcounter{enumi}{31}
\tightlist
\item
  设 B 是拓扑空间，\((\mathrm{E},p)\) 是 B-空间。设 \((\mathrm{V}_{i})_{i\in\mathrm{I}}\) 是 B 的可数化覆盖，并且对每个 \(i\in I\)，\(V_{i}\)-空间 \(E_{V_{i}}\) 具有截面延拓性质。
\end{enumerate}

\begin{enumerate}
\def\labelenumi{\alph{enumi})}
\item
  我们假设 B 上存在连续的局部有限单位分解 \((f_{i})\)，使得对所有 i 都有 \(V_{i} = f_{i}^{-1}(]0,1])\)。令 \(g: B \to [0,1]\) 为连续映射。置 \(A = g^{-1}(1)\)，\(W = g^{-1}(]0,1])\)。令 s 为 E 在 W 上的截面。证明存在 E 的一个截面，使其在 A 上与 s 重合。（考虑如下集合中的一个极大元：其元素是二元组 \((J,t)\)，其中 J 是 I 的子集，t 是 E 在 \(W \cup \bigcup_{i \in J} V_{i}\) 上的截面，且在 A 上与 s 重合；该集合赋予如下序关系：\((J,t) \prec (J',t')\) 当 \(J \subset J'\) 且 \(t'\) 延拓 t。）
\item
  证明 B-空间 E 具有截面延拓性质。
\end{enumerate}

